\documentclass[10pt,a4paper]{amsart}
\usepackage{amsmath,amssymb,amsthm,graphicx}
\usepackage{verbatim}
\newenvironment{reference}{\comment}{\endcomment}
%\newenvironment{reference}{}{}
\newenvironment{slogan}{\comment}{\endcomment}
\newenvironment{history}{\comment}{\endcomment}

% For commutative diagrams we use Xy-pic
\usepackage[all]{xy}

% We use 2cell for 2-commutative diagrams.
\xyoption{2cell}
% Theorem environments.
%
\theoremstyle{plain}
\newtheorem{theorem}[subsection]{Theorem}
\newtheorem{proposition}[subsection]{Proposition}
\newtheorem{lemma}[subsection]{Lemma}

\theoremstyle{definition}
\newtheorem{definition}[subsection]{Definition}
\newtheorem{example}[subsection]{Example}
\newtheorem{exercise}[subsection]{Exercise}
\newtheorem{situation}[subsection]{Situation}

\theoremstyle{remark}
\newtheorem{remark}[subsection]{Remark}
\newtheorem{remarks}[subsection]{Remarks}

\numberwithin{equation}{subsection}

% Macros
%
\def\lim{\mathop{\mathrm{lim}}\nolimits}
\def\colim{\mathop{\mathrm{colim}}\nolimits}
\def\Spec{\mathop{\mathrm{Spec}}}
\def\Hom{\mathop{\mathrm{Hom}}\nolimits}
\def\Ext{\mathop{\mathrm{Ext}}\nolimits}
\def\SheafHom{\mathop{\mathcal{H}\!\mathit{om}}\nolimits}
\def\SheafExt{\mathop{\mathcal{E}\!\mathit{xt}}\nolimits}
\def\Sch{\mathit{Sch}}
\def\Mor{\mathop{\mathrm{Mor}}\nolimits}
\def\Ob{\mathop{\mathrm{Ob}}\nolimits}
\def\Sh{\mathop{\mathit{Sh}}\nolimits}
\def\NL{\mathop{N\!L}\nolimits}
\def\CH{\mathop{\mathrm{CH}}\nolimits}
\def\proetale{{pro\text{-}\acute{e}tale}}
\def\etale{{\acute{e}tale}}
\def\QCoh{\mathit{QCoh}}
\def\Ker{\mathop{\mathrm{Ker}}}
\def\Im{\mathop{\mathrm{Im}}}
\def\Coker{\mathop{\mathrm{Coker}}}
\def\Coim{\mathop{\mathrm{Coim}}}

% Boxtimes
%
\DeclareMathSymbol{\boxtimes}{\mathbin}{AMSa}{"02}

%
% Macros for moduli stacks/spaces
%
\def\QCohstack{\mathcal{QC}\!\mathit{oh}}
\def\Cohstack{\mathcal{C}\!\mathit{oh}}
\def\Spacesstack{\mathcal{S}\!\mathit{paces}}
\def\Quotfunctor{\mathrm{Quot}}
\def\Hilbfunctor{\mathrm{Hilb}}
\def\Curvesstack{\mathcal{C}\!\mathit{urves}}
\def\Polarizedstack{\mathcal{P}\!\mathit{olarized}}
\def\Complexesstack{\mathcal{C}\!\mathit{omplexes}}
% \Pic is the operator that assigns to X its picard group, usage \Pic(X)
% \Picardstack_{X/B} denotes the Picard stack of X over B
% \Picardfunctor_{X/B} denotes the Picard functor of X over B
\def\Pic{\mathop{\mathrm{Pic}}\nolimits}
\def\Picardstack{\mathcal{P}\!\mathit{ic}}
\def\Picardfunctor{\mathrm{Pic}}
\def\Deformationcategory{\mathcal{D}\!\mathit{ef}}
\usepackage{cleveref}
\crefname{originallemma}{original lemma}{original lemmas}
\Crefname{originallemma}{Original lemma}{Original lemmas}
\makeatletter
\newlabel{derived-lemma-negative-exts}{{06XS}{source}{}{sourcepointer-1}{}}
\newlabel{derived-lemma-negative-exts@cref}{{[originallemma][0][]06XS}{[1][0][]source}{}{}{}}
\makeatother
\crefname{originalreference}{original reference}{original references}
\Crefname{originalreference}{Original reference}{Original references}
\makeatletter
\newlabel{lemma-compute-ext-perfect}{{08ID}{source}{}{sourcepointer-2}{}}
\newlabel{lemma-compute-ext-perfect@cref}{{[originalreference][0][]08ID}{[1][0][]source}{}{}{}}
\makeatother
\crefname{originalreference}{original reference}{original references}
\Crefname{originalreference}{Original reference}{Original references}
\makeatletter
\newlabel{lemma-identify-pseudo-coherent-noetherian}{{08E8}{source}{}{sourcepointer-3}{}}
\newlabel{lemma-identify-pseudo-coherent-noetherian@cref}{{[originalreference][0][]08E8}{[1][0][]source}{}{}{}}
\makeatother
\crefname{originalreference}{original reference}{original references}
\Crefname{originalreference}{Original reference}{Original references}
\makeatletter
\newlabel{lemma-pseudo-coherent-hocolim}{{0DJN}{source}{}{sourcepointer-4}{}}
\newlabel{lemma-pseudo-coherent-hocolim@cref}{{[originalreference][0][]0DJN}{[1][0][]source}{}{}{}}
\makeatother
\crefname{originalreference}{original reference}{original references}
\Crefname{originalreference}{Original reference}{Original references}
\makeatletter
\newlabel{lemma-tensor-perfect}{{09J5}{source}{}{sourcepointer-5}{}}
\newlabel{lemma-tensor-perfect@cref}{{[originalreference][0][]09J5}{[1][0][]source}{}{}{}}
\makeatother
\crefname{originallemma}{original lemma}{original lemmas}
\Crefname{originallemma}{Original lemma}{Original lemmas}
\makeatletter
\newlabel{sites-cohomology-lemma-dual-perfect-complex}{{08JJ}{source}{}{sourcepointer-6}{}}
\newlabel{sites-cohomology-lemma-dual-perfect-complex@cref}{{[originallemma][0][]08JJ}{[1][0][]source}{}{}{}}
\makeatother
\crefname{originalreference}{original reference}{original references}
\Crefname{originalreference}{Original reference}{Original references}
\makeatletter
\newlabel{theorem-approximation}{{08ES}{source}{}{sourcepointer-7}{}}
\newlabel{theorem-approximation@cref}{{[originalreference][0][]08ES}{[1][0][]source}{}{}{}}
\makeatother
\title{Computing Ext against a bounded flat complex with proper support by a perfect object and by a pseudo-coherent dual}
\author{The Stacks Project Authors}
\date{stacks-project:spaces-perfect.tex:local-snapshot-20261012; GNU FDL 1.2 or later}
\begin{document}
\maketitle
\noindent\textit{Editorial note.} All delivered mathematical text is the author's own, copied from the pinned Stacks Project chapter file spaces-perfect.tex; no statement or proof step was rewritten, reordered or removed.\par\medskip
\noindent\textit{Editorial note.} The counted proof is the author's complete argument: an approximation-and-cone reduction for part (A) and a homotopy-colimit-of-duals construction with truncation bookkeeping for part (B), both carried out inside the delivered span\par\medskip
\noindent\textit{Editorial note.} The label lemma-compute-ext names the counted statement, and the proof environment immediately following it in the source is delivered verbatim as its proof.\par\medskip
\noindent\textit{Editorial note.} The counted statement is carried by the author's own lemma environment in the chapter source.\par\medskip
\noindent\textit{Allowed standard prerequisites (editorial).}\begin{itemize}
\item The approximation theorem for objects of the derived category by perfect complexes, cited from the same chapter
\item Ext vanishing for cones with truncated cohomology support, and the presentation of pseudo-coherent objects as homotopy colimits of perfect objects
\item Duality for perfect complexes on ringed sites and preservation of perfectness under derived tensor and pushforward
\end{itemize}\medskip
\ifdefined\hypertarget\hypertarget{sourcepointer-1}{}\fi
\noindent\textit{Source reference.} Original lemma 06XS: sources/stacks-chow/tags.txt, source lines 8861--8861. Prerequisite pointer invoked at proof line(s) 6074; it names a standard fact of a sibling chapter that the counted argument uses, and every step of the counted proof built on it is retained verbatim in the delivered body.\par\medskip
\ifdefined\hypertarget\hypertarget{sourcepointer-2}{}\fi
\noindent\textit{Source reference.} Original reference 08ID: sources/stacks-spaces-perfect/pinned-original-source.tex, source lines 5946--5946. Prerequisite pointer invoked at proof line(s) 6076; it names a standard fact of the same chapter that the counted argument uses, and every step of the counted proof built on it is retained verbatim in the delivered body.\par\medskip
\ifdefined\hypertarget\hypertarget{sourcepointer-3}{}\fi
\noindent\textit{Source reference.} Original reference 08E8: sources/stacks-spaces-perfect/pinned-original-source.tex, source lines 3324--3324. Prerequisite pointer invoked at proof line(s) 6083; it names a standard fact of the same chapter that the counted argument uses, and every step of the counted proof built on it is retained verbatim in the delivered body.\par\medskip
\ifdefined\hypertarget\hypertarget{sourcepointer-4}{}\fi
\noindent\textit{Source reference.} Original reference 0DJN: sources/stacks-spaces-perfect/pinned-original-source.tex, source lines 4537--4537. Prerequisite pointer invoked at proof line(s) 6084; it names a standard fact of the same chapter that the counted argument uses, and every step of the counted proof built on it is retained verbatim in the delivered body.\par\medskip
\ifdefined\hypertarget\hypertarget{sourcepointer-5}{}\fi
\noindent\textit{Source reference.} Original reference 09J5: sources/stacks-spaces-perfect/pinned-original-source.tex, source lines 5794--5794. Prerequisite pointer invoked at proof line(s) 6095; it names a standard fact of the same chapter that the counted argument uses, and every step of the counted proof built on it is retained verbatim in the delivered body.\par\medskip
\ifdefined\hypertarget\hypertarget{sourcepointer-6}{}\fi
\noindent\textit{Source reference.} Original lemma 08JJ: sources/stacks-chow/tags.txt, source lines 10852--10852. Prerequisite pointer invoked at proof line(s) 6088; it names a standard fact of a sibling chapter that the counted argument uses, and every step of the counted proof built on it is retained verbatim in the delivered body.\par\medskip
\ifdefined\hypertarget\hypertarget{sourcepointer-7}{}\fi
\noindent\textit{Source reference.} Original reference 08ES: sources/stacks-spaces-perfect/pinned-original-source.tex, source lines 3786--3786. Prerequisite pointer invoked at proof line(s) 6053; it names a standard fact of the same chapter that the counted argument uses, and every step of the counted proof built on it is retained verbatim in the delivered body.\par\medskip
\begin{lemma}
\label{lemma-compute-ext}
Let $S$ be a scheme. Let $f : X \to B$ be a morphism of algebraic spaces
over $S$, $E \in D(\mathcal{O}_X)$, and $\mathcal{F}^\bullet$ a complex
of $\mathcal{O}_X$-modules. Assume
\begin{enumerate}
\item $B$ is Noetherian,
\item $f$ is locally of finite type and quasi-separated,
\item $E \in D^-_{\textit{Coh}}(\mathcal{O}_X)$,
\item $\mathcal{G}^\bullet$ is a bounded complex of coherent
$\mathcal{O}_X$-module flat over $B$ with support proper over $B$.
\end{enumerate}
Then the following two statements are true
\begin{enumerate}
\item[(A)] for every $m \in \mathbf{Z}$ there exists a perfect object $K$
of $D(\mathcal{O}_B)$ and functorial maps
$$
\alpha^i_\mathcal{F} :
\Ext^i_{\mathcal{O}_X}(E,
\mathcal{G}^\bullet \otimes_{\mathcal{O}_X} f^*\mathcal{F})
\longrightarrow
H^i(B, K \otimes^\mathbf{L}_{\mathcal{O}_B} \mathcal{F})
$$
for $\mathcal{F}$ quasi-coherent on $B$
compatible with boundary maps (see proof)
such that $\alpha^i_\mathcal{F}$ is an isomorphism for $i \leq m$, and
\item[(B)] there exists a pseudo-coherent $L \in D(\mathcal{O}_B)$
and functorial isomorphisms
$$
\Ext^i_{\mathcal{O}_B}(L, \mathcal{F}) \longrightarrow
\Ext^i_{\mathcal{O}_X}(E,
\mathcal{G}^\bullet \otimes_{\mathcal{O}_X} f^*\mathcal{F})
$$
for $\mathcal{F}$ quasi-coherent on $B$ compatible with boundary maps.
\end{enumerate}
\end{lemma}

\begin{proof}
Proof of (A). Suppose $\mathcal{G}^i$ is nonzero only for $i \in [a, b]$.
We may replace $X$ by a quasi-compact open neighbourhood of the union
of the supports of $\mathcal{G}^i$. Hence we may assume $X$ is Noetherian.
In this case $X$ and $f$ are quasi-compact and quasi-separated.
Choose an approximation $P \to E$ by a perfect complex $P$ of
$(X, E, -m - 1 + a)$
(possible by Theorem \ref{theorem-approximation}).
Then the induced map
$$
\Ext^i_{\mathcal{O}_X}(E,
\mathcal{G}^\bullet \otimes_{\mathcal{O}_X} f^*\mathcal{F})
\longrightarrow
\Ext^i_{\mathcal{O}_X}(P,
\mathcal{G}^\bullet \otimes_{\mathcal{O}_X} f^*\mathcal{F})
$$
is an isomorphism for $i \leq m$. Namely, the kernel, resp.\ cokernel of this
map is a quotient, resp.\ submodule of
$$
\Ext^i_{\mathcal{O}_X}(C,
\mathcal{G}^\bullet \otimes_{\mathcal{O}_X} f^*\mathcal{F})
\quad\text{resp.}\quad
\Ext^{i + 1}_{\mathcal{O}_X}(C,
\mathcal{G}^\bullet \otimes_{\mathcal{O}_X} f^*\mathcal{F})
$$
where $C$ is the cone of $P \to E$. Since $C$ has vanishing cohomology
sheaves in degrees $\geq -m - 1 + a$ these $\Ext$-groups are zero
for $i \leq m + 1$ by
Derived Categories, Lemma \ref{derived-lemma-negative-exts}.
This reduces us to the case that
$E$ is a perfect complex which is Lemma \ref{lemma-compute-ext-perfect}.
The statement on boundaries is explained in the proof of
Lemma \ref{lemma-compute-ext-perfect}.

\medskip\noindent
Proof of (B). As in the proof of (A) we may assume $X$ is Noetherian.
Observe that $E$ is pseudo-coherent by
Lemma \ref{lemma-identify-pseudo-coherent-noetherian}.
By Lemma \ref{lemma-pseudo-coherent-hocolim} we can write
$E = \text{hocolim} E_n$ with $E_n$ perfect and $E_n \to E$ inducing
an isomorphism on truncations $\tau_{\geq -n}$. Let $E_n^\vee$
be the dual perfect complex
(Cohomology on Sites, Lemma \ref{sites-cohomology-lemma-dual-perfect-complex}).
We obtain an inverse system $\ldots \to E_3^\vee \to E_2^\vee \to E_1^\vee$
of perfect objects. This in turn gives rise to an inverse system
$$
\ldots \to K_3 \to K_2 \to K_1\quad\text{with}\quad
K_n = Rf_*(E_n^\vee \otimes_{\mathcal{O}_X}^\mathbf{L} \mathcal{G}^\bullet)
$$
perfect on $Y$, see Lemma \ref{lemma-tensor-perfect}.
By Lemma \ref{lemma-compute-ext-perfect} and its proof and
by the arguments in the previous paragraph (with $P = E_n$)
for any quasi-coherent $\mathcal{F}$ on $Y$ we have
functorial canonical maps
$$
\xymatrix{
& \Ext^i_{\mathcal{O}_X}(E,
\mathcal{G}^\bullet \otimes_{\mathcal{O}_X} f^*\mathcal{F})
\ar[ld] \ar[rd] \\
H^i(Y, K_{n + 1} \otimes_{\mathcal{O}_Y}^\mathbf{L} \mathcal{F})
\ar[rr] & &
H^i(Y, K_n \otimes_{\mathcal{O}_Y}^\mathbf{L} \mathcal{F})
}
$$
which are isomorphisms for $i \leq n + a$.
Let $L_n = K_n^\vee$ be the dual perfect complex.
Then we see that $L_1 \to L_2 \to L_3 \to \ldots$
is a system of perfect objects in $D(\mathcal{O}_Y)$
such that for any quasi-coherent $\mathcal{F}$ on $Y$
the maps
$$
\Ext^i_{\mathcal{O}_Y}(L_{n + 1}, \mathcal{F})
\longrightarrow
\Ext^i_{\mathcal{O}_Y}(L_n, \mathcal{F})
$$
are isomorphisms for $i \leq n + a - 1$. This implies that
$L_n \to L_{n + 1}$ induces an isomorphism on truncations
$\tau_{\geq -n - a + 2}$ (hint: take cone of $L_n \to L_{n + 1}$
and look at its last nonvanishing cohomology sheaf).
Thus $L = \text{hocolim} L_n$ is pseudo-coherent, see
Lemma \ref{lemma-pseudo-coherent-hocolim}. The mapping property
of homotopy colimits gives that
$\Ext^i_{\mathcal{O}_Y}(L, \mathcal{F}) =
\Ext^i_{\mathcal{O}_Y}(L_n, \mathcal{F})$
for $i \leq n + a - 3$ which finishes the proof.
\end{proof}

\end{document}
